\documentclass{article}
\usepackage{graphicx} % Required for inserting images
\usepackage{amsmath}
\usepackage{amssymb}
\usepackage[colorlinks=true, linkcolor=blue, urlcolor=blue]{hyperref}

\DeclareMathOperator*{\argmax}{arg\,max}
\DeclareMathOperator*{\argmin}{arg\,min}
\title{Annealed RL}
\date{December 2025}

\begin{document}

\maketitle

\section{Introduction}

When post-training LLMs we often have distinct phases. 

\subsection{SFT}
Offline RL in the form of SFT which is a set $D$ of question answer pairs, $D = \{ (q,a) \mid a \sim \pi_{\text{ref}}(\cdot \mid q)$ where $\pi_{\text{ref}}$ is some unknown distribution that we ideally want to match. In practice this can be answers written by human annotations or high quality question-answer completions scrapped off the internet. For our purposes we treat an answer as a single sample from this high quality, reference distribution. In theory if we wanted to match an alternative distribution such as a LLM like GPT-5, we could $n$ samples from it and thus that would be $a_1, \ldots, a_n \sim \pi_{\text{ref}} = \pi_{\text{GPT-5}}$. The point is $\pi_{\text{ref}}$ is some distribution which are samples are coming from and we want to maximize our probability of sampling from this distribution. We would like to match our model distribution $\pi_\theta(\cdot \mid q)$ to a
reference distribution $\pi_{\text{ref}}(\cdot \mid q)$ by minimizing the KL
divergence
\begin{equation}
D_{\mathrm{KL}}\!\left(
\pi_{\text{ref}}(\cdot \mid q)
\;\middle\|\;
\pi_\theta(\cdot \mid q)
\right)
=
\mathbb{E}_{a \sim \pi_{\text{ref}}(\cdot \mid q)}
\left[
\log \pi_{\text{ref}}(a \mid q)
-
\log \pi_\theta(a \mid q)
\right].
\end{equation}

Taking the expectation over questions $q \sim D$ yields the optimization
objective
\begin{equation}
\min_\theta \;
\mathbb{E}_{q \sim D}
\mathbb{E}_{a \sim \pi_{\text{ref}}(\cdot \mid q)}
\left[
\log \pi_{\text{ref}}(a \mid q)
-
\log \pi_\theta(a \mid q)
\right].
\end{equation}

The first term $\log \pi_{\text{ref}}(a \mid q)$ does not depend on $\theta$ and
can therefore be dropped from the optimization. This gives
\begin{equation}
\theta
=
\argmin_\theta
\;
\mathbb{E}_{q \sim D}
\mathbb{E}_{a \sim \pi_{\text{ref}}(\cdot \mid q)}
\left[
- \log \pi_\theta(a \mid q)
\right],
\end{equation}
which is the standard supervised fine-tuning (SFT) objective.

\subsection{RLVR}

A newer phase is the RLVR idea where you can perform online, on-policy improvements for question answer pairs that have a correct answer, e.g math, code, etc. In this case we focus on single-turn RLVR which means we only have one chat-turn. 
Usually you perform SFT on some basic question answer pairs first to be able to get some reward and then you perform online RL sampling questions and completing them with the model that is being optimized. That is you now have your dataset $D = \{ q \}$ be questions and you sample $G$ completions from your model and optimize over it. 
Then the optimization term is some variant of policy-gradient

\begin{align}
\theta &= \argmin_\theta \; \mathbb{E}_{q \sim D} \mathbb{E}_{a \sim \pi_\theta(\cdot \mid q)}
\left[
-A(a) \log \pi_\theta(a \mid q)
\right] \\
&\approx \argmin_\theta \; \mathbb{E}_{q \sim D} \left[ \frac{1}{G}\sum_{g = 1}^G
\left[
-A(a) \log \pi_\theta(a_g \mid q)
\right] \right]\\
\end{align}

\section{Main Idea}

The main idea is to understand what we mean by $\pi_\theta(a \mid q)$ for a question answer $(q,a)$ pair. Since LLMs are auto-regressive we obtain $\log \pi_\theta ( a \mid q) = \sum_{i=1}^T \log \pi_\theta(a^t \mid q, a^{<t})$ where $a^{< t}$ is the concatenation of all tokens before the $t^{\text{th}}$ token. Extending this notation let $a^{i:j}$ be the concatenation of tokens from $(a^i, a^{i+1}, \ldots, a^{j-1}, a^j)$ for $i \leq j$. 

The main question we are trying to answer then is: \textit{Is ideal to go from an sft objective where your answer $a^{0:T}$ is generated from $\pi_{\text{ref}}$ and then to an RLVR objective where $a^{0:T}$ is generated from $\pi_\theta$? Can you interpolate through both by annealing from the SFT objective to the RLVR objective through the sampling tokens?} 

The idea is to make $T$ tokens as a function of time instead of switching after the phases are done. An example could be for the first 10 steps you do full SFT on 8192 tokens and then for the next 20 steps you do SFT on 5000 tokens and the rest $\sim 3000$ tokens are generated on policy from your policy that you are optimizing.

More formally let $T(s): [0, \text{max\_steps}] \rightarrow [0, T]$  be a monotonically decreasing, non-negative function that takes an input step and returns back some integer in the range of 0 to max tokens. Then we define the training objective on the $s$ step as 

\begin{equation} \label{mainObj}
\begin{aligned}
- \mathbb{E}_{q \sim D} \Big[
&\mathbb{E}_{a^{0:T(s)} \sim \pi_{\text{ref}}(\cdot \mid q)}
\log \pi_\theta(a^{0 : T(s)} \mid q) \\
&+ \mathbb{E}_{a^{T(s): T} \sim \pi_\theta(\cdot \mid q)}
A(a^{T(s): T}) \log \pi_\theta(a^{T(s): T} \mid q)
\Big]
\end{aligned}
\end{equation}

\section{Experiment ideas}


There are two ways to think of the idea as of now. One is from a SFT and RLVR mix perspective and the other is to just focus on the RL term.

Starting with the first we want to train on examples with the joint objective from \ref{mainObj}. This would involed us selecting some pair of sft dataset that we can RL on like \href{https://huggingface.co/datasets/AI-MO/NuminaMath-CoT}{NuminaMath}. The issue with this is we have to show that this is better than pure RL and SFT since if we have the completions why not just use pure sft then pure RLVR why the interpolation. 

The next perspective is to think of this as an introduction o RL where we only train on the RL part of the objective in essense just mask out the SFT loss. This can be seen as a better approach as we are annealing RL and is not a direct comparison to SFT. 

To do experiments we would ideally choose Qwen 8B and then do it on math very similar to \href{https://www.notion.so/sagnikm/Who-is-Adam-SGD-Might-Be-All-We-Need-For-RLVR-In-LLMs-1cd2c74770c080de9cbbf74db14286b6}{Who is Adam?} blogpost. 

We would also have to probably ablate over which schedules we want 

\section{Things to do}

A few things are needed to be done for this. The first is we need a JAX native inference setup. We need to figure out fast datasets and environments. Then we need a trainer that is very flexible to do these experiments. 

\pagebreak 
\bibliographystyle{plain}   % or unsrt, alpha, abbrv, etc.
\bibliography{ref}
\textbf{TODO} 

\end{document}
